\documentclass[12pt,a4paper]{article}

% Codificação e idioma
\usepackage[utf8]{inputenc}
\usepackage[portuguese]{babel}

% Pacotes úteis
\usepackage{booktabs}
\usepackage{longtable}
\usepackage{graphicx}
\usepackage{hyperref}
\usepackage{geometry}
\usepackage{indentfirst}
\usepackage{float}

% Margens
\geometry{margin=2.5cm}


\newcommand{\projetotitulo}{Introdução à Manufatura Aditiva}
\newcommand{\grupoid}{Grupo Manufatura Aditiva}
\newcommand{\sprintnum}{5}
\newcommand{\sprintperiodo}{22/06 a 05/07/2026}
\newcommand{\scrummaster}{Lucas Moura}
\newcommand{\gitrepo}{https://github.com/viniciusnovelo/Introducao-a-Manufatura-Aditiva}
\newcommand{\gitcommit}{} % hash do commit desta entrega

\begin{document}

\begin{center}
{\Large\bfseries Relatório de Sprint \sprintnum}\\
{\large \projetotitulo}\\
\grupoid{} \ | \ \sprintperiodo{} \ | \ SM: \scrummaster
\end{center}

\section{Sprint Backlog}
\footnote{hash: \gitcommit}

\begin{longtable}{|c|p{3.8cm}|p{1.9cm}|p{2.2cm}|p{1.7cm}|p{4.2cm}|}
\toprule
ID & História & Estimativa & Responsável & Status & Critério de Aceitação \\
\midrule
1 & Criação do Instagram para Divulgação & Sprint 1 & Lucas Moura & Concluído & Perfil criado, configurado com identidade visual do projeto e contendo ao menos uma publicação inicial. \\ \hline

2 & Criação do formulário de divulgação & Sprint 1 & Vinícius & Concluído & Formulário criado, testado e apto para receber inscrições dos participantes. \\ \hline

3 & Início da divulgação para comunidade interna e externa & Sprint 1 & Vinícius & Concluído & Divulgação realizada em ao menos uma escola ou meio de comunicação definido pela equipe. \\ \hline

4 & Criação de slides para as aulas de SolidEdge para os calouros & Sprint 1 & Paulo & Concluído & Slides elaborados e revisados pela equipe antes da aplicação das aulas. \\ \hline

5 & Criação do quadro Kanban & Sprint 1 & Lucas Leão & Concluído & Quadro criado contendo colunas de acompanhamento e tarefas cadastradas. \\ \hline

6 & Coleta e organização de dados das inscrições dos alunos & Sprint 1 & Paulo & Concluído & Dados armazenados e organizados em planilha. \\ \hline

7 & Agendamento das salas e laboratórios que serão utilizados & Sprint 1 & Leonardo & Concluído & Reservas confirmadas junto à instituição para todas as aulas previstas. \\ \hline

8 & Verificação e manutenção das impressoras 3D e softwares no laboratório & Sprint 1 & Yuri & Concluído & Equipamentos testados e softwares funcionando corretamente antes do início do curso. \\ \hline

9 & Criação de relatório de Sprint 1 & Sprint 1 & Lucas Leão & Concluído & Relatório elaborado contendo atividades realizadas, resultados e pendências. \\ \hline

10 & Criação do repositório no Github & Sprint 1 & Vinícius & Concluído & Repositório criado com acesso para a equipe e estrutura inicial organizada. \\ \hline

11 & Treinamento da equipe nas ferramentas de modelagem 3D & Sprint 1 e 2 & Lucas Moura & Concluído & Equipe capacitada para executar operações básicas de modelagem no software. \\ \hline

12 & Elaboração dos materiais para as aulas 1 e 2 do curso & Sprint 2 & Equipe & Concluído & Materiais preparados e revisados antes da realização das aulas. \\ \hline

13 & Realização da aula 1 & Sprint 2 & Equipe & Concluído & Aula ministrada conforme cronograma e com registro de presença dos alunos. \\ \hline

14 & Realização da aula 2 & Sprint 2 & Equipe & Concluído & Aula ministrada conforme cronograma e com registro de presença dos alunos. \\ \hline

15 & Criação de relatório de Sprint 2 & Sprint 2 & Leonardo, Paulo, Yuri & Concluído & Relatório elaborado contendo atividades realizadas, resultados e pendências. \\ \hline

16 & Atualização do quadro Kanban & Sprint 2 & Lucas Leão & Concluído & Todas as tarefas da sprint atualizadas conforme o andamento das atividades. \\ \hline

17 & Atualização do repositório Github & Sprint 2 & Vinícius & Concluído & Materiais e documentos da sprint enviados e organizados no repositório. \\ \hline

18 & Treinamento da equipe nas ferramentas de impressão 3D & Sprint 3 & Lucas Moura & Concluído & Equipe apta a operar as impressoras e configurar impressões básicas. \\ \hline

19 & Elaboração dos materiais para as aulas 3 e 4 do curso & Sprint 3 & Equipe & Concluído & Materiais preparados e aprovados pela equipe antes das aulas. \\ \hline

20 & Realização da aula 3 & Sprint 3 & Equipe & Concluído & Aula ministrada conforme cronograma e com registro de presença dos alunos. \\ \hline

21 & Criação de relatório de Sprint 3 & Sprint 3 & Leonardo, Paulo, Yuri & Concluído & Relatório elaborado contendo atividades realizadas, resultados e pendências. \\ \hline

22 & Atualização do quadro Kanban & Sprint 3 & Lucas Leão & Concluído & Todas as tarefas da sprint atualizadas conforme o andamento das atividades. \\ \hline

23 & Atualização do repositório Github & Sprint 3 & Vinícius & Concluído & Materiais e documentos da sprint enviados e organizados no repositório. \\ \hline

24 & Elaboração dos materiais para as aulas 5 e 6 do curso & Sprint 4 & Lucas Moura & Concluído & Materiais finalizados, revisados e disponíveis para aplicação das aulas. \\ \hline

25 & Realização da aula 4 & Sprint 4 & Equipe & Concluído & Aula ministrada conforme cronograma e com registro de presença dos alunos. \\ \hline

26 & Realização da aula 5 & Sprint 4 & Equipe & Concluído & Aula ministrada conforme cronograma e com registro de presença dos alunos. \\ \hline

27 & Criação de relatório de Sprint 4 & Sprint 4 & Leonardo, Paulo, Yuri & Concluído & Relatório elaborado contendo atividades realizadas, resultados e pendências. \\ \hline

28 & Atualização do quadro Kanban & Sprint 4 & Lucas Leão & Concluído & Todas as tarefas da sprint atualizadas conforme o andamento das atividades. \\ \hline

29 & Atualização do repositório Github & Sprint 4 & Vinícius & Concluído & Materiais e documentos da sprint enviados e organizados no repositório. \\ \hline

30 & Realização da aula 6 & Sprint 5 & Equipe & Concluído & Aula final realizada e conteúdo previsto concluído. \\ \hline

31 & Criação de formulário para coleta de feedback & Sprint 6 & Leonardo, Yuri & Em andamento & Formulário criado contendo perguntas objetivas e discursivas sobre o curso. \\ \hline

32 & Coleta de feedback dos alunos através de formulário & Sprint 6 & Leonardo, Yuri & Em andamento & Respostas coletadas de parte significativa dos participantes do curso. \\ \hline

33 & Criação de relatório de Sprint 5 & Sprint 5 & Leonardo, Paulo, Yuri & Concluído & Relatório elaborado contendo atividades realizadas, resultados e pendências. \\ \hline

34 & Atualização do quadro Kanban & Sprint 6 & Lucas Leão & Em andamento & Todas as tarefas da sprint atualizadas conforme o andamento das atividades. \\ \hline

35 & Atualização do repositório Github & Sprint 6 & Vinícius & Em andamento & Materiais e documentos da sprint enviados e organizados no repositório. \\ \hline

36 & Criação do relatório final & Sprint 6 & A definir & Não iniciado & Relatório final concluído contendo descrição completa do projeto, resultados e conclusões. \\

\bottomrule
\end{longtable}

\section{Burndown Chart}
\begin{figure}[h!]
\centering
\includegraphics[width=1.0\textwidth]{Imagens/burndown_sprint\sprintnum.png}}
\caption{Burndown -- Sprint \sprintnum}
\end{figure}

% Análise do burndown (mínimo 3 linhas):
O gráfico acima representa o engajamento e o trabalho efetivo realizado pelo grupo durante a Sprint 5, sendo que a linha azul representa o progresso ideal no decorrer de cada dia e a linha laranja o progresso realizado efetivamente. Nesta Sprint, as atividades propostas foram realizadas dentro do prazo previsto.

\section{Artefatos Produzidos}

Durante a Sprint 5 foi produzido e confeccionado uma peça modelo para última aula, além de várias impressões que foram entregues aos alunos como lembrancinhas ao final da última aula.

\begin{figure}[H]
\centering

\begin{minipage}{0.48\textwidth}
    \centering
    \includegraphics[width=\textwidth]{Imagens/Aula 6 - 1.jpeg}
    \caption{Aula 27/06}
\end{minipage}
\hfill
\begin{minipage}{0.48\textwidth}
    \centering
    \includegraphics[width=\textwidth]{Imagens/Aula 6 - 2.jpeg}
        \caption{Aula 27/06}
\end{minipage}

\vspace{2.0cm}

\begin{minipage}{0.48\textwidth}
    \centering
    \includegraphics[width=\textwidth]{Imagens/Aula 6 - 3.jpeg}
        \caption{Aula 27/06}
\end{minipage}
\hfill
\begin{minipage}{0.48\textwidth}
    \centering
    \includegraphics[width=\textwidth]{Imagens/Peça Modelo.jpeg}
        \caption{Peça Modelo}
\end{minipage}
\end{figure}

\begin{figure}[H]
\centering

\begin{minipage}{0.41\textwidth}
    \centering
    \includegraphics[width=\textwidth]{Imagens/Lembrancinha.jpeg}
    \caption{Lembrancinha}
\end{minipage}
\hfill
\begin{minipage}{0.55\textwidth}
    \centering
    \includegraphics[width=\textwidth]{Imagens/Registro final.jpeg}
        \caption{Registro Final}
\end{minipage}

\end{figure}


\section{Atualização com o Parceiro}
Na Sprint 5, os parceiros participaram da última aula e finalizaram o aprendizado de todo o conteúdo proposto, desde modelagem até a configuração, parametrização e por fim, impressão 3D. Muitos dos parceiros não concluíram o curso com os requisitos propostos, entretando, o último encontro contou com a participação de 90\% dos alunos realmente engajados e com maior participação nos encontros.  

\section{Participação do Calouro}
Na Sprint 5, os calouros auxiliaram nas configurações da impressora 3D e na apresentação da última aula, além de contribuirem na organização e definição dos alunos que concluiram o curso com sucesso.

\section{Impedimentos e Desvios}
Na preparação para a última aula, houve dificuldade na configuração e manuseio da impressora 3D devido à pontos de degradação do equipamento, entretato, os ajustes necessários foram realizados e não houve nenhum imprevisto ou impedimento na última aula.

\section{Retrospectiva}
\subsection*{O que funcionou bem}
As ações previstas para esta sprint foi foram devidamente finalizadas com os conteúdos apresentados no último encontro e com a entrega dos certificados aos alunos. 

\subsection*{O que precisa melhorar}
Tendo em vista os próximos passos para conclusão do projeto, é necessário maior atenção para os prazos e entregas finais e individuais. 

\subsection*{Ação concreta para a próxima sprint}
Visando estes pontos, o grupo se reunirá novamente e será feito um acompanhamento coletivo de todos os membros para garantir a qualidade e pontualidade das entregas finais. 

\section{Planejamento da Próxima Sprint}
Para próxima Sprint, será levantado o feedback dos alunos que receberam os certificados a fim de avaliar a qualidade do curso ofertado e possíveis críticas ou pontos de melhoria. Ademais, tanto o repositório e o quadro Kanban serão finalizados e avaliados pelos demais membros do grupo e os membros se organizarão para apresentação final do projeto.
\end{document}